\chapter{OneLake and Microsoft Fabric Foundations}
\index{Microsoft Fabric}\index{OneLake}\index{Lakehouse}\index{shortcut}\index{Delta Lake}\index{Power BI}\index{capacity units}%
\index{medallion architecture}\index{data virtualization}\index{Fabric capacity}%

\begin{objectives}
\begin{itemize}
    \item Explain Microsoft Fabric as a unified SaaS analytics platform and distinguish its major workloads.
    \item Reason about Fabric capacities, F SKUs, capacity units, workspaces, domains, and tenant-level organization.
    \item Define OneLake as a single logical data lake per tenant and describe how workspaces, folders, and lakehouses project into it.
    \item Describe Delta-Parquet as the native open storage pattern for reliable lakehouse tables.
    \item Create and evaluate OneLake shortcuts for internal and external data virtualization scenarios.
    \item Build a first lakehouse lab that reads shortcut data with PySpark without copying the underlying files.
    \item Architect a multi-cloud retail data virtualization pattern that unifies AWS S3 history and Dataverse CRM data through OneLake.
\end{itemize}
\end{objectives}

\begin{crosscloud}
Although this book teaches \textbf{Microsoft Fabric}, OneLake is built directly on the Azure column below, and the Week~1 storage foundations map almost one-to-one across the other clouds.

\vspace{2mm}
{\footnotesize
\begin{tabular}{@{}p{2.8cm}p{3.0cm}p{3.0cm}p{3.0cm}@{}}
\toprule
\textbf{Capability} & \textbf{AWS} & \textbf{Azure} & \textbf{GCP} \\
\midrule
Object storage & Amazon S3 & ADLS Gen2 / Blob & Cloud Storage \\
Hot tier & S3 Standard & Hot & Standard \\
Infrequent access & S3 Standard-IA & Cool / Cold & Nearline \\
Archive (instant) & Glacier Instant Retrieval & Cold & Coldline \\
Deep archive & Glacier Deep Archive & Archive & Archive \\
Lifecycle rules & Lifecycle policies & Lifecycle management & Object Lifecycle Mgmt \\
Versioning & S3 Versioning & Blob versioning & Object Versioning \\
Immutability (WORM) & Object Lock & Immutable blob storage & Bucket Lock \\
Cross-region copy & CRR / SRR & Object replication (GRS) & Dual/multi-region + Transfer \\
Online transfer & DataSync & AzCopy & Storage Transfer Service \\
Offline transfer & Snow Family & Data Box & Transfer Appliance \\
\addlinespace
Design durability & 11 nines & 11 nines (LRS) / 16 (GRS) & 11 nines \\
Availability SLA & 99.9\% (Standard) & 99.9\% (RA-GRS read 99.99\%) & 99.95\% (multi-region) \\
Encryption at rest & SSE-S3 / SSE-KMS / SSE-C & SSE + CMK via Key Vault & Google-managed / CMEK / CSEK \\
Access control & IAM \& bucket policies, ACLs & Azure RBAC, SAS, POSIX ACLs & IAM, ACLs, signed URLs \\
Pricing levers & Storage + requests + egress + retrieval & Storage + operations + egress + retrieval & Storage + operations + egress + retrieval \\
Consistency & Strong read-after-write & Strong read-after-write & Strong (global) \\
\bottomrule
\end{tabular}}

\vspace{1mm}
The same storage layer reappears as a managed abstraction across the stack: \textbf{Fabric OneLake} is a tenant-wide lake over ADLS Gen2 (Delta--Parquet), \textbf{Snowflake} persists data as immutable micro-partitions on the provider's object store, and \textbf{MongoDB}'s WiredTiger is a node-local engine rather than object storage---the contrast is itself instructive.
\end{crosscloud}

\begin{keyterms}
\textbf{Fabric capacity}: a pool of compute resources, measured through capacity units, that powers Fabric workloads. \quad
\textbf{OneLake}: the tenant-wide logical data lake for Microsoft Fabric. \quad
\textbf{Lakehouse}: a Fabric item that combines data lake storage, Delta tables, SQL access, and Spark access. \quad
\textbf{Shortcut}: a OneLake pointer to data stored elsewhere, exposed as if it were local. \quad
\textbf{Delta-Parquet}: Parquet files plus Delta Lake transaction logs that support ACID table semantics.
\end{keyterms}

\section{Week 1 Roadmap}
This chapter is the first week of the Microsoft Fabric Master Data Engineer program. Monday introduces the platform. Tuesday establishes OneLake and the lakehouse storage contract. Wednesday focuses on shortcuts and data virtualization. Thursday is a guided project in which you create a workspace, provision a lakehouse, create a shortcut to public blob data, and read the data from a notebook. Friday closes with a global retailer architecture pattern that uses OneLake as the unifying namespace for multi-cloud analytics.

The week deliberately begins with foundations rather than ingestion tools. In Fabric, storage architecture choices influence every later decision: Spark jobs write Delta tables, warehouses can query lakehouse data through SQL endpoints, Power BI can use Direct Lake over curated models, domains organize data products, and governance metadata depends on workspace and item boundaries. Microsoft describes Fabric as an end-to-end analytics platform that brings data engineering, data integration, warehousing, real-time intelligence, data science, and business intelligence into a single SaaS experience \cite{microsoftFabricOverview}. OneLake is the storage foundation of that experience \cite{microsoftOneLakeOverview}.

\begin{notebox}
This chapter intentionally uses the term \emph{data product}. In Fabric, a useful data product is not only a table. It is a governed bundle of storage, transformations, ownership, quality expectations, lineage, security, and consumption paths.
\end{notebox}

\section{Monday: Introduction to Microsoft Fabric}
\index{Microsoft Fabric!overview}\index{software as a service}\index{workloads}
Microsoft Fabric is best understood as a unified SaaS analytics operating environment rather than a single data tool. It combines experiences that historically lived in separate products: Data Factory for integration, Synapse-style data engineering and warehousing, Real-Time Intelligence, Data Science, and Power BI. The unification matters because the platform gives these workloads shared identity, shared capacity, shared governance, shared workspace organization, and shared OneLake storage.

\begin{definitionbox}
\textbf{Microsoft Fabric} is a unified SaaS analytics platform that integrates data movement, data engineering, data warehousing, real-time analytics, data science, governance, and Power BI over a common tenant, workspace, capacity, and OneLake foundation.
\end{definitionbox}

A traditional Azure analytics estate often contains independent services: storage accounts, Spark pools, SQL pools, Data Factory instances, Key Vaults, Power BI workspaces, monitoring resources, and security assignments. Fabric does not remove the need for architecture. It changes the default unit of architecture. Workspaces and Fabric items become the visible work surface; capacities become the shared compute boundary; OneLake becomes the default data lake; and Power BI semantic delivery is integrated instead of bolted on at the end.

\subsection{The Workload Map}
Fabric workloads are specialized experiences over a shared platform. A data engineer should know their purpose and their storage relationship.

\begin{table}[H]
\centering
\small
\begin{tabular}{@{}p{3.2cm}p{5.6cm}p{5.2cm}@{}}
\toprule
\textbf{Workload} & \textbf{Primary role} & \textbf{Storage and serving pattern} \\
\midrule
Data Engineering & Build lakehouses, notebooks, Spark jobs, and Delta tables & Writes files and tables in OneLake, commonly in medallion layers \\
Data Factory & Ingest, transform, and orchestrate data pipelines & Moves or transforms data into lakehouses, warehouses, or external systems \\
Data Warehouse & Serve relational analytics with SQL & Uses warehouse storage and can integrate with lakehouse data through SQL experiences \\
Real-Time Intelligence & Process event streams, logs, and telemetry & Eventstreams and KQL databases support streaming and operational analytics \\
Data Science & Explore, train, track, and operationalize models & Uses notebooks, experiments, models, and lakehouse data \\
Power BI & Model and visualize business data & Consumes lakehouse and warehouse data through import, DirectQuery, or Direct Lake \\
\bottomrule
\end{tabular}
\caption{Major Fabric workloads and their relationship to the shared analytics platform.}
\label{tab:chapter1-workloads}
\end{table}

\begin{notebox}
Do not treat workloads as independent silos. The same customer table may be landed by Data Factory, cleaned by Spark, served by a warehouse, certified in a semantic model, secured by domain governance, and visualized in Power BI. Fabric's value is the shared operating model across those steps.
\end{notebox}

\subsection{Capacities, F SKUs, and Capacity Units}
\index{capacity}\index{F SKU}\index{capacity unit}
A Fabric capacity is the compute pool assigned to one or more workspaces. Microsoft licenses Fabric through capacity SKUs and concepts that determine how much work the platform can execute \cite{microsoftFabricCapacity}. F SKUs represent Fabric capacities with different amounts of available compute. Capacity units, often abbreviated as CUs, are the practical abstraction used to reason about consumption across workloads.

\begin{definitionbox}
A \textbf{capacity unit} is an abstract measure of Fabric compute throughput. It lets the platform meter diverse work---Spark sessions, warehouse queries, Data Factory activities, Power BI interactions, and real-time workloads---against a shared capacity pool.
\end{definitionbox}

Capacity planning is not only a procurement exercise. It is an engineering feedback loop. The same capacity can be affected by long Spark notebooks, expensive SQL queries, concurrent refreshes, pipeline bursts, and report traffic. A lakehouse design that creates many small files may slow downstream reads and burn capacity. A semantic model that unnecessarily imports data may duplicate storage and refresh work. A pipeline that copies external data that could have been virtualized through a shortcut may increase both latency and compute cost.

\begin{tipbox}
Start capacity design by classifying workloads by shape: interactive reports, scheduled batch jobs, ad hoc notebooks, streaming ingestion, and administrative operations. Then test concurrency, not just single-job performance.
\end{tipbox}

\subsection{Workspaces and Domains}
\index{workspace}\index{domain}
Workspaces are collaboration and item containers. They hold lakehouses, warehouses, notebooks, pipelines, semantic models, reports, dataflows, eventstreams, KQL databases, and related artifacts. A workspace is also the practical boundary for assigning capacity and permissions. Domains provide a higher-level governance organization, allowing related workspaces and data products to be grouped by business area such as Finance, Retail Operations, Supply Chain, or Customer Experience \cite{microsoftDomains}.

A master data engineer treats workspace design as information architecture. A single shared workspace for every team becomes ungovernable. A workspace per developer fragments lineage and reuse. A useful middle ground aligns workspaces to data products, environments, or bounded domains. For example, \texttt{retail-sales-dev}, \texttt{retail-sales-test}, and \texttt{retail-sales-prod} can live under a Retail domain, while customer support analytics lives in a Customer domain.

\begin{pitfall}
\textbf{Using workspaces as folders.} Workspaces carry permissions, capacity assignment, lineage, deployment, and operational meaning. If you create workspaces only to make a navigation tree look tidy, you will eventually fight inconsistent access and duplicated data products.
\end{pitfall}

\subsection{Monday Architecture Lens}
The platform view produces three questions that recur throughout the book. First, what is the compute boundary? In Fabric, the answer is normally a capacity assigned to workspaces. Second, what is the collaboration and security boundary? The answer is normally a workspace, refined by item permissions and downstream semantic model security. Third, what is the data boundary? The answer is OneLake plus the source systems linked through shortcuts or pipelines. Good Fabric architecture keeps those three boundaries aligned enough to operate, but not so tightly coupled that every change requires a tenant redesign.

\section{Tuesday: OneLake Concepts}
\index{OneLake!concepts}\index{OneDrive for Data}\index{Delta-Parquet}
Microsoft presents OneLake as the ``OneDrive for Data'' for an organization \cite{microsoftOneLakeOverview}. The phrase is useful because it emphasizes a single logical namespace across the tenant. Users do not need to choose a storage account every time they create a lakehouse. Fabric items place data into OneLake paths managed by the platform, and authorized tools can access those files through Fabric experiences and supported APIs.

\begin{definitionbox}
\textbf{OneLake} is the single logical data lake for a Fabric tenant. It stores Fabric data items and exposes a unified namespace so workloads can share data without every team provisioning separate storage accounts.
\end{definitionbox}

The single-lake idea does not mean every team sees every file. Authorization still matters. Workspaces, item permissions, data access roles, sharing, endorsement, and governance controls determine who can discover and use data. The value of OneLake is that storage becomes a shared substrate rather than a collection of disconnected lakes.

\subsection{Workspaces, Items, Folders, and Paths}
A Fabric lakehouse item has a storage layout that separates files and tables. The \texttt{Files} area is for arbitrary folders and files, including raw landing data. The \texttt{Tables} area is for managed Delta tables that Fabric can expose through Spark, SQL analytics endpoints, and Power BI-oriented experiences. This distinction becomes central when you design medallion layers.

\begin{notebox}
A folder of Parquet files in \texttt{Files} is not automatically the same thing as a managed lakehouse table. A Delta table has a transaction log and table metadata that enable reliable reads, writes, schema evolution, and query interoperability.
\end{notebox}

\subsection{Delta-Parquet as the Native Open Format}
\index{Delta Lake}\index{Parquet}\index{ACID}
Fabric lakehouses use open data formats. Parquet provides columnar file storage with efficient compression and predicate pushdown \cite{apacheParquet}. Delta Lake adds a transaction log over Parquet files to provide ACID table behavior on object storage \cite{armbrust2020delta,deltaLakeDocs}. The lakehouse architecture pattern unifies warehouse-style analytics and data-science workloads over open storage \cite{armbrust2021lakehouse}. Fabric adopts this direction by making Delta tables a first-class lakehouse construct \cite{microsoftFabricLakehouse}.

\begin{definitionbox}
\textbf{Delta-Parquet} means table data is stored as Parquet files and table state is managed by a Delta transaction log. The combination supports atomic commits, time-aware metadata, schema enforcement, and consistent reads by multiple engines.
\end{definitionbox}

Delta matters because distributed engines do not write one neat database file. Spark jobs write many files, and concurrent jobs can overlap. Without a transaction log, readers may see partial results, duplicate files, or inconsistent schemas. With Delta, a table commit is a versioned metadata operation over a set of Parquet files.

\subsection{Medallion Thinking in OneLake}
\index{medallion architecture}
The medallion architecture organizes data into Bronze, Silver, and Gold layers. Bronze stores raw or lightly standardized data. Silver stores cleaned, conformed, deduplicated, and enriched data. Gold stores business-ready aggregates and dimensional models for consumption \cite{databricksMedallion}. OneLake does not force this layout, but Fabric lakehouses make it natural to model those layers through lakehouses, folders, schemas, or tables.

\begin{figure}[H]
    \centering
    \resizebox{\textwidth}{!}{%
    \begin{tikzpicture}[
        zone/.style={rectangle, rounded corners, draw=primary, thick, fill=primary!5, align=center, minimum width=2.7cm, minimum height=1.1cm, font=\small\bfseries},
        medal/.style={cylinder, shape border rotate=90, aspect=0.25, draw=primary, thick, fill=backcolour, align=center, minimum width=2.5cm, minimum height=1.4cm, font=\small\bfseries},
        arrow/.style={draw, thick, -Latex, draw=primary}
    ]
        \node[zone] (src) at (0, 0) {Sources\\ERP, CRM, S3, IoT};
        \node[zone] (shortcut) at (3.5, 0) {OneLake\\Shortcuts};
        \node[medal, fill=brown!12] (bronze) at (7.0, 0) {Bronze\\Raw Delta};
        \node[medal, fill=gray!15] (silver) at (10.5, 0) {Silver\\Conformed};
        \node[medal, fill=yellow!18] (gold) at (14.0, 0) {Gold\\Business Data};
        \node[zone] (serve) at (17.5, 0) {Direct Lake\\Power BI};

        \draw[arrow] (src) -- node[above, font=\scriptsize]{virtualize} (shortcut);
        \draw[arrow] (shortcut) -- node[above, font=\scriptsize]{read} (bronze);
        \draw[arrow] (bronze) -- node[above, font=\scriptsize]{clean} (silver);
        \draw[arrow] (silver) -- node[above, font=\scriptsize]{model} (gold);
        \draw[arrow] (gold) -- node[above, font=\scriptsize]{serve} (serve);

        \node[note] at (8.8, -1.45) {Open Delta-Parquet tables in OneLake preserve interoperability across Spark, SQL, and BI workloads.};
    \end{tikzpicture}%
    }
    \caption{A OneLake-centered medallion pattern for Microsoft Fabric.}
    \label{fig:chapter1-onelake-medallion}
\end{figure}

\subsection{Tuesday Engineering Lens}
OneLake encourages teams to separate logical data access from physical storage administration. That separation is powerful, but it can hide operational reality. A table still has files. Files still have sizes, partitions, formats, schemas, and update patterns. A shortcut still depends on a remote system. A report still depends on a semantic model and refresh behavior. The mature Fabric engineer uses OneLake to simplify access while continuing to reason about storage contracts.

\section{Wednesday: OneLake Shortcuts}
\index{OneLake!shortcuts}\index{data virtualization}\index{ADLS Gen2}\index{Amazon S3}\index{Dataverse}
A shortcut is a pointer inside OneLake to data stored somewhere else. Instead of copying files from an external source into OneLake, Fabric can expose the remote data through the OneLake namespace. Microsoft documents shortcuts for sources such as ADLS Gen2, Amazon S3, Dataverse, and other Fabric locations \cite{microsoftOneLakeShortcuts}. For a data engineer, shortcuts are one of Fabric's most important design primitives because they separate logical access from physical movement.

\begin{definitionbox}
A \textbf{OneLake shortcut} is a virtual folder reference that appears inside OneLake while the underlying files remain in their source location. The shortcut lets Fabric workloads read data through the OneLake namespace without first copying it.
\end{definitionbox}

\subsection{Internal and External Shortcuts}
Internal shortcuts point to data already in OneLake, often in another workspace or item. They support reuse without duplicating files. External shortcuts point to data outside OneLake, such as ADLS Gen2, Amazon S3, or Dataverse. Both patterns can simplify architecture, but they have different governance and operational implications.

\begin{table}[H]
\centering
\small
\begin{tabular}{@{}p{3.2cm}p{4.8cm}p{5.4cm}@{}}
\toprule
\textbf{Shortcut type} & \textbf{Typical use} & \textbf{Engineering questions} \\
\midrule
Internal OneLake & Share curated data between workspaces without copying & Who owns the source? What happens if the source table schema changes? How is lineage communicated? \\
ADLS Gen2 & Reuse Azure lake data during migration or coexistence & Which identity accesses storage? Are firewall, region, and folder conventions compatible? \\
Amazon S3 & Virtualize historical or partner data in AWS & What are latency, egress, credentials, and consistency implications? \\
Dataverse & Bring business application data into analytical pipelines & Which tables are needed? How are privacy, row-level access, and refresh expectations handled? \\
\bottomrule
\end{tabular}
\caption{Shortcut types and architecture questions.}
\label{tab:chapter1-shortcut-types}
\end{table}

\subsection{When Virtualization Beats Copying}
Shortcut virtualization is valuable when data is large, shared, governed by another platform, updated by another team, or too expensive to duplicate. It is also useful during migrations: a team can build Fabric transformation logic against external paths, then later decide whether selected data should be physically materialized into OneLake Delta tables.

\begin{tipbox}
Use shortcuts for access and exploration; use managed Delta tables for reliable curated products. A shortcut to raw external data is not a substitute for Silver and Gold tables with clear ownership, quality checks, and contracts.
\end{tipbox}

\subsection{Risks and Controls}
Shortcuts do not remove physics. Remote data can be slower than local data. External cloud access can involve network latency, availability dependencies, credential management, and possible data-transfer costs. External schemas can change without Fabric knowing in advance. The source system remains the system of record, so deletion or restructuring at the source may break downstream Fabric jobs.

\begin{pitfall}
\textbf{Treating virtual data as governed data.} A shortcut makes external data visible through OneLake, but it does not automatically certify quality, freeze schema, apply business definitions, or guarantee service-level agreements. Curate important shortcut data into managed Delta tables before exposing it broadly.
\end{pitfall}

\subsection{Shortcut Decision Matrix}
Shortcut decisions should be documented explicitly. The following matrix is a practical review aid.

\begin{table}[H]
\centering
\small
\begin{tabular}{@{}p{4.0cm}p{4.5cm}p{4.5cm}@{}}
\toprule
\textbf{Question} & \textbf{Shortcut-friendly answer} & \textbf{Materialize-first answer} \\
\midrule
How large is the source? & Too large or costly to copy before proving value & Small enough to copy cheaply and repeatedly \\
Who owns source quality? & Another team owns it and publishes a clear contract & Your team must repair quality before any consumption \\
How stable is schema? & Stable or source publishes change notifications & Volatile and likely to break downstream code \\
What is the consumption path? & Exploration, profiling, or controlled engineering jobs & Certified BI, finance reporting, or audited outputs \\
What is the latency target? & Batch or exploratory access tolerates remote reads & Interactive dashboards require predictable local performance \\
\bottomrule
\end{tabular}
\caption{Shortcut versus materialization decision matrix.}
\label{tab:chapter1-shortcut-matrix}
\end{table}

\section{Thursday: Hands-on Project}
\index{hands-on lab}\index{workspace!creation}\index{lakehouse!creation}\index{PySpark}
The project for Week 1 creates the minimum useful Fabric data engineering environment: a workspace, a lakehouse, a shortcut, and a notebook that reads the shortcut. The scenario assumes you have access to a Fabric tenant, an available capacity or trial capacity, and permissions to create a workspace and lakehouse.

\subsection{Goal and Architecture}
You will create a workspace named \texttt{fabric-de-week01}, add a lakehouse named \texttt{lh\_onelake\_foundations}, create a shortcut named \texttt{public\_samples}, and read files exposed by that shortcut with PySpark. The shortcut target can be any public Azure Blob container that your tenant can access. If your instructor provides a specific URL, use that URL. If not, use a public sample container approved for your training environment.

\begin{notebox}
Fabric user interface labels change over time. When a button label differs slightly, follow the action intent: create a workspace, create a lakehouse item, open the lakehouse explorer, create a shortcut under \texttt{Files}, select the Azure Blob or ADLS-compatible source, and validate that files appear without uploading them.
\end{notebox}

\subsection{Step-by-Step Actions}
\begin{enumerate}
    \item Open the Microsoft Fabric portal and confirm that the Fabric experience is enabled for your tenant.
    \item Select \textbf{Workspaces}, then create a new workspace named \texttt{fabric-de-week01}. Assign it to a Fabric capacity if your organization requires explicit capacity selection.
    \item Inside the workspace, choose \textbf{New item}, then create a \textbf{Lakehouse} named \texttt{lh\_onelake\_foundations}.
    \item In the lakehouse explorer, expand \texttt{Files}. Choose the command to create a new shortcut.
    \item Select the external source type for the public sample container. For an Azure Blob public container, choose the source option your tenant exposes for Blob or ADLS-compatible access.
    \item Enter shortcut name \texttt{public\_samples}. Enter the account, container, and folder information supplied by your instructor or the public training dataset.
    \item Complete the connection step. If authentication is requested for a public source, choose the anonymous or organizational method that matches the published source and your tenant policy.
    \item Confirm that \texttt{Files/public\_samples} appears in the lakehouse and that its folders or files are visible. You have virtualized data; you have not copied it into the lakehouse.
    \item Create a new notebook attached to the lakehouse. Name it \texttt{nb\_week01\_shortcut\_read}.
    \item Run the PySpark listing below, adjusting the relative path and file format to match the public dataset.
\end{enumerate}

\begin{lstlisting}[language=Python,caption={PySpark notebook snippet for reading data exposed through a OneLake shortcut.}]
from pyspark.sql import functions as F

# In a Fabric notebook attached to lh_onelake_foundations, Files paths are
# reachable relative to the lakehouse mount. Adjust the child folder and file
# extension to match the public sample container used by your instructor.
shortcut_path = "Files/public_samples/sales/*.parquet"

sales_df = spark.read.parquet(shortcut_path)

profile_df = (
    sales_df
    .select(
        F.count("*").alias("row_count"),
        F.countDistinct("CustomerId").alias("customer_count"),
        F.min("OrderDate").alias("first_order_date"),
        F.max("OrderDate").alias("last_order_date"),
        F.sum("SalesAmount").alias("total_sales_amount"),
    )
)

profile_df.show(truncate=False)

# Persist a curated Delta table in the lakehouse. The source files remain in
# the external container; this table is a managed Fabric lakehouse asset.
clean_df = (
    sales_df
    .withColumn("OrderDate", F.to_date("OrderDate"))
    .withColumn("SalesAmount", F.col("SalesAmount").cast("decimal(18,2)"))
    .dropDuplicates(["OrderId"])
)

clean_df.write.format("delta").mode("overwrite").saveAsTable("silver_sales_orders")
\end{lstlisting}

\subsection{Representing the Shortcut as Configuration}
In production, connection details should be managed through deployment processes, environment-specific parameters, and secure credentials. The following JSON is not a secret file. It is a conceptual deployment record that captures the intent of a shortcut so it can be reviewed consistently.

\begin{lstlisting}[language=json,caption={Conceptual JSON record for documenting a OneLake shortcut target.}]
{
  "workspace": "fabric-de-week01",
  "lakehouse": "lh_onelake_foundations",
  "shortcut": {
    "name": "public_samples",
    "location": "Files/public_samples",
    "sourceType": "AzureBlobPublicContainer",
    "sourceUri": "https://examplepublicdata.blob.core.windows.net/sales",
    "copyData": false,
    "intendedLayer": "Bronze exploration",
    "owner": "data-engineering-training",
    "qualityContract": "not certified; curate before BI consumption"
  }
}
\end{lstlisting}

\subsection{Optional SQL Check}
After the notebook writes a managed Delta table, use the SQL analytics endpoint or a SQL query cell to check record counts and date ranges. This listing is intentionally simple: Week 1 validates access and table creation, not warehouse performance tuning.

\begin{lstlisting}[language=SQL,caption={SQL validation after writing a managed Delta table from shortcut data.}]
SELECT
    COUNT(*) AS row_count,
    COUNT(DISTINCT CustomerId) AS customer_count,
    MIN(OrderDate) AS first_order_date,
    MAX(OrderDate) AS last_order_date,
    SUM(SalesAmount) AS total_sales_amount
FROM silver_sales_orders;
\end{lstlisting}

\begin{tipbox}
After a successful shortcut read, immediately write down three facts: the owner of the external data, the expected refresh or update pattern, and the downstream table that will become the certified copy. These facts prevent shortcut sprawl.
\end{tipbox}

\subsection{Validation Checklist}
A successful lab has observable evidence:
\begin{itemize}
    \item The workspace exists and is assigned to the expected capacity.
    \item The lakehouse appears in the workspace and opens without errors.
    \item The shortcut folder appears under \texttt{Files} and exposes external files.
    \item The notebook reads the shortcut path without uploading files manually.
    \item A managed Delta table such as \texttt{silver\_sales\_orders} is created from the shortcut data.
    \item You can explain which data remains external and which data is now managed in the lakehouse.
\end{itemize}

\section{Friday: Use Case and Architecture}
\index{multi-cloud architecture}\index{retail analytics}\index{Dataverse}\index{Amazon S3}
A global retailer has ten years of historical sales data stored in Amazon S3 by region. Current customer and loyalty activity lives in Dataverse because front-office teams use Dynamics-based applications. Executives want a single Power BI view of customer lifetime value, store performance, and campaign effectiveness. The data engineering team wants to avoid a large one-time migration before proving value.

\subsection{Architecture Strategy}
The design uses OneLake shortcuts to virtualize both historical S3 data and current Dataverse data into a Fabric lakehouse. Bronze reads remain virtualized at first. Silver tables standardize customer, order, product, calendar, and region fields into managed Delta tables. Gold tables publish conformed metrics for Power BI through a semantic model using Direct Lake where appropriate \cite{microsoftDirectLake}.

\begin{figure}[H]
    \centering
    \resizebox{\textwidth}{!}{%
    \begin{tikzpicture}[
        source/.style={rectangle, rounded corners, draw=codegray, thick, fill=white, align=center, minimum width=2.6cm, minimum height=1.1cm, font=\small\bfseries},
        virtual/.style={rectangle, rounded corners, draw=primary, thick, fill=primary!8, align=center, minimum width=3.0cm, minimum height=1.1cm, font=\small\bfseries},
        lake/.style={cylinder, shape border rotate=90, aspect=0.25, draw=primary, thick, fill=backcolour, align=center, minimum width=2.6cm, minimum height=1.5cm, font=\small\bfseries},
        consume/.style={rectangle, rounded corners, draw=codegreen!60!black, thick, fill=green!7, align=center, minimum width=2.7cm, minimum height=1.1cm, font=\small\bfseries},
        arrow/.style={draw, thick, -Latex, draw=primary}
    ]
        \node[source] (s3) at (0, 1.4) {AWS S3\\Historical Sales};
        \node[source] (dv) at (0, -1.4) {Dataverse\\CRM + Loyalty};
        \node[virtual] (short) at (4.0, 0) {OneLake Shortcuts\\No Initial Copy};
        \node[lake, fill=brown!12] (bronze) at (8.0, 0) {Bronze\\Virtual Raw};
        \node[lake, fill=gray!15] (silver) at (11.7, 0) {Silver\\Conformed Delta};
        \node[lake, fill=yellow!18] (gold) at (15.4, 0) {Gold\\Retail Metrics};
        \node[consume] (pbi) at (19.0, 0.9) {Power BI\\Direct Lake};
        \node[consume] (wh) at (19.0, -0.9) {Warehouse\\SQL Analytics};

        \draw[arrow] (s3) -- (short);
        \draw[arrow] (dv) -- (short);
        \draw[arrow] (short) -- (bronze);
        \draw[arrow] (bronze) -- node[above, font=\scriptsize]{standardize} (silver);
        \draw[arrow] (silver) -- node[above, font=\scriptsize]{aggregate} (gold);
        \draw[arrow] (gold) -- (pbi);
        \draw[arrow] (gold) -- (wh);

        \node[note, text width=13cm, align=center] at (9.8, -2.2) {Governance binds the design: domain ownership, source contracts, lineage, sensitivity labels, quality checks, and certified semantic models.};
    \end{tikzpicture}%
    }
    \caption{Multi-cloud retail virtualization through OneLake shortcuts and medallion curation.}
    \label{fig:chapter1-retail-virtualization}
\end{figure}

\subsection{Design Decisions}
The first decision is to avoid copying all S3 history on day one. A shortcut makes the data visible to Fabric jobs, which lets engineers profile schema, quality, and partitioning before committing to a migration. The second decision is to materialize curated Silver and Gold tables in OneLake. Business users should not query raw external folders directly. The third decision is to represent Dataverse as a governed source with privacy-aware curation: customer identifiers, consent fields, and loyalty attributes require explicit access rules.

\begin{table}[H]
\centering
\small
\begin{tabular}{@{}p{3.5cm}p{5.0cm}p{5.0cm}@{}}
\toprule
\textbf{Concern} & \textbf{Shortcut-first decision} & \textbf{Curation decision} \\
\midrule
Historical sales & Use S3 shortcut for profiling and incremental value & Convert validated partitions into Silver Delta tables \\
Current CRM & Use Dataverse shortcut for governed analytical access & Mask or restrict sensitive fields before Gold publication \\
Cost and latency & Avoid bulk migration before value is proven & Optimize high-value queries into local Delta and semantic models \\
Ownership & Source teams retain operational systems & Domain data owners certify curated tables and metrics \\
BI consumption & Do not point reports at raw virtual folders & Publish Gold metrics through Direct Lake semantic models \\
\bottomrule
\end{tabular}
\caption{Retail architecture decisions for shortcut-first modernization.}
\label{tab:chapter1-retail-decisions}
\end{table}

\begin{pitfall}
\textbf{Skipping the Silver contract.} Multi-cloud virtualization can make data look unified before it is semantically unified. Customer keys, time zones, currencies, returns, store hierarchies, and consent fields must be standardized before executives trust a global report.
\end{pitfall}

\section{Operational Guidance for Week 1}
A Week 1 Fabric environment can already fail in production-like ways. A workspace can be created in the wrong capacity. A shortcut can point to the wrong folder. A notebook can read external files slowly because partitioning is poor. A user can assume a shortcut means the data has been certified. The data engineer's job is to make these assumptions visible.

\subsection{Governance Checklist}
\begin{itemize}
    \item Assign every workspace to a domain or document why it is temporary.
    \item Name lakehouses, shortcuts, notebooks, and tables consistently.
    \item Document source ownership and update frequency for every shortcut.
    \item Separate raw shortcut exploration from certified Silver and Gold tables.
    \item Record sensitivity and privacy expectations before publishing semantic models.
    \item Test whether intended users can access the data and whether unintended users cannot.
\end{itemize}

\subsection{Quality and Observability}
Even in the first lab, build the habit of recording quality observations. Count rows, inspect null rates, check duplicate keys, record source file paths, and save the notebook output that proves the shortcut worked. In later chapters these checks become formal data quality gates. Week 1 is where the discipline begins.

\begin{lstlisting}[language=Python,caption={Small data quality profile for a newly read shortcut data frame.}]
from pyspark.sql import functions as F

required_columns = ["OrderId", "CustomerId", "OrderDate", "SalesAmount"]
missing_columns = sorted(set(required_columns) - set(sales_df.columns))
if missing_columns:
    raise ValueError(f"Missing required columns: {missing_columns}")

quality_df = sales_df.select(
    F.count("*").alias("rows"),
    F.sum(F.col("OrderId").isNull().cast("int")).alias("null_order_ids"),
    F.sum(F.col("CustomerId").isNull().cast("int")).alias("null_customer_ids"),
    F.countDistinct("OrderId").alias("distinct_orders"),
)

quality_df.show(truncate=False)
\end{lstlisting}

\subsection{Interview Q\&A}
\begin{enumerate}
    \item \textbf{Q: Why is OneLake described as OneDrive for Data?}\\
    A: Because it provides a single logical tenant-wide namespace where Fabric data items can be discovered and accessed through shared platform controls, rather than forcing every team to manage separate storage accounts.
    \item \textbf{Q: When would you use a shortcut instead of copying data?}\\
    A: Use a shortcut when the source is large, externally owned, frequently reused, or still being evaluated. Copy or materialize when you need curated quality, predictable performance, stable schemas, and certified consumption.
    \item \textbf{Q: What is the relationship between Parquet and Delta Lake?}\\
    A: Parquet stores columnar files; Delta Lake adds the transaction log and metadata that make those files behave like reliable tables for concurrent engines.
    \item \textbf{Q: Why do capacities matter to data engineers?}\\
    A: Capacity is the shared compute pool. Inefficient Spark jobs, warehouse queries, refreshes, and report usage compete for it, so engineering choices affect performance and cost.
    \item \textbf{Q: Why should a BI report not point directly to a raw external shortcut?}\\
    A: Raw shortcut data usually lacks semantic standardization, quality checks, privacy treatment, and performance guarantees. A report should consume certified Gold tables or semantic models.
    \item \textbf{Q: How do domains change the operating model?}\\
    A: Domains group related workspaces and data products under business ownership. They make governance and discovery align with organizational accountability rather than with accidental workspace lists.
\end{enumerate}

\section{Further Reading}
Begin with Microsoft's overview of Fabric \cite{microsoftFabricOverview}, OneLake \cite{microsoftOneLakeOverview}, OneLake shortcuts \cite{microsoftOneLakeShortcuts}, lakehouses \cite{microsoftFabricLakehouse}, capacities and licenses \cite{microsoftFabricCapacity}, domains \cite{microsoftDomains}, and Direct Lake \cite{microsoftDirectLake}. For storage foundations, read the Delta Lake paper \cite{armbrust2020delta}, the lakehouse architecture paper \cite{armbrust2021lakehouse}, Apache Parquet documentation \cite{apacheParquet}, and Delta Lake documentation \cite{deltaLakeDocs}. For medallion design vocabulary, compare Fabric implementations with the medallion architecture description \cite{databricksMedallion}.

\section*{Key Takeaways}
\begin{itemize}
    \item Fabric is a unified SaaS analytics platform, not merely a collection of renamed Azure services.
    \item Capacities and capacity units turn compute into a shared engineering constraint across Spark, SQL, pipelines, real-time workloads, and Power BI.
    \item OneLake provides a single logical tenant-wide data lake, while workspaces and items provide practical collaboration and security boundaries.
    \item Fabric lakehouses use open Delta-Parquet patterns so Spark, SQL, and BI workloads can interoperate over reliable tables.
    \item Shortcuts virtualize data from internal OneLake locations and external systems such as ADLS Gen2, Amazon S3, and Dataverse.
    \item Virtualization avoids unnecessary copying, but curated Silver and Gold Delta tables are still required for trusted business consumption.
    \item A sound Week 1 architecture already includes ownership, naming, quality, privacy, capacity, and consumption decisions.
\end{itemize}

\section{Exercises}
\begin{enumerate}
    \item \textbf{(Conceptual)} Explain the difference between a Fabric workspace, a domain, a capacity, and a lakehouse. Give one governance decision associated with each.
    \item \textbf{(Conceptual)} Why is a shortcut not the same thing as a certified table? Describe at least three risks that remain after a shortcut is created.
    \item \textbf{(Hands-on)} Create a development workspace and lakehouse. Add a shortcut under \texttt{Files}, then take a screenshot or write a short observation proving that the data was not uploaded manually.
    \item \textbf{(Hands-on)} Modify the PySpark listing to read CSV or JSON data from your shortcut, infer or provide a schema, and write a managed Delta table named \texttt{bronze\_source\_profile}.
    \item \textbf{(Design)} Design a workspace naming convention for development, test, and production Fabric environments in a retail organization. Include capacity assignment and domain placement.
    \item \textbf{(Design)} For the global retailer, decide which data should remain virtualized and which should be materialized into OneLake. Justify your decision using cost, latency, ownership, and BI trust.
    \item \textbf{(Interview)} In five sentences, defend Delta-Parquet as the native lakehouse format for Fabric data engineering.
    \item \textbf{(Operational)} Create a Week 1 runbook entry that documents shortcut owner, source URI, authentication method, expected update frequency, downstream tables, and escalation contact.
\end{enumerate}

\section{Capstone Project: A Multi-Cloud Retail Landing Zone with OneLake Shortcuts}\label{sec:ch1-capstone}
\index{capstone project}\index{landing zone}\index{multi-cloud retail}\index{OneLake!shortcuts}\index{Silver Delta table}

\begin{capstonebox}
\textbf{Mission.} Design and implement a Week~1 retail landing zone that uses Microsoft Fabric as the SaaS control plane, a trial-backed workspace as the collaboration boundary, OneLake as the analytical namespace, and shortcuts as the virtualization layer. Your goal is to let analysts query historical sales in Amazon S3 and current CRM data from Dataverse without first performing a bulk migration. You will then write one curated Silver Delta table that becomes the trusted starting point for downstream BI.

\textbf{Estimated time.} 3--5 hours, depending on tenant permissions, source availability, and whether your class uses real S3 and Dataverse sources or instructor-provided public Blob and simulated CRM files.

\textbf{What you'll practice.}
\begin{itemize}
    \item Creating a Fabric workspace on trial capacity and assigning work to a Lakehouse item.
    \item Using OneLake shortcuts to expose external data in place instead of copying raw files.
    \item Organizing a Bronze zone that separates virtual raw folders from managed curated tables.
    \item Reading shortcut paths with PySpark and producing a unified Delta-Parquet Silver table.
    \item Explaining the difference between virtualized source access and certified analytical data products.
\end{itemize}
\end{capstonebox}

\subsection*{Scenario}
A global retailer operates stores and e-commerce channels in North America, Europe, and Asia-Pacific. Ten years of historical order, return, and store sales files are stored by region in AWS S3 because the original point-of-sale modernization program ran on AWS. Current customer profile, loyalty status, and campaign consent data lives in Dataverse because the customer service and marketing teams use Dynamics-based applications.

The executive request is simple: create a single Fabric landing zone that can answer customer lifetime value, regional sales, and loyalty conversion questions without waiting for a large data migration. The engineering constraint is equally important: historical S3 files and operational CRM exports must remain in their source systems during Week~1. Fabric should provide unified access, profiling, and curation while avoiding unnecessary duplication of raw data. This is the Friday architecture scenario made concrete as a hands-on capstone.

\subsection*{Requirements}
Build the smallest production-shaped solution that proves the architecture:
\begin{itemize}
    \item A Fabric workspace named \texttt{fabric-de-week01-retail} assigned to a Fabric trial capacity or another instructor-approved capacity \cite{microsoftFabricTrial}.
    \item A Lakehouse named \texttt{lh\_retail\_landing} in that workspace.
    \item A Bronze folder structure under \texttt{Files/bronze} with shortcut locations for historical sales and CRM data.
    \item OneLake shortcuts to an Amazon S3 source for historical sales, or to an instructor-approved public Blob source that simulates S3 if S3 credentials are unavailable.
    \item A OneLake shortcut to Dataverse CRM data, or to an instructor-approved simulated CRM folder if Dataverse permissions are unavailable.
    \item A PySpark notebook named \texttt{nb\_capstone\_retail\_silver} attached to the Lakehouse.
    \item Notebook logic that reads both shortcut sources in place, standardizes schemas, joins or unions data as appropriate, and writes a managed Silver Delta table named \texttt{silver\_retail\_customer\_sales}.
    \item No bulk copy of raw S3, Blob, or Dataverse data into OneLake. Only the curated Silver Delta output is materialized in the Lakehouse.
\end{itemize}

\begin{notebox}
If your tenant does not allow S3 or Dataverse shortcuts, use public Blob or instructor-provided files as stand-ins. Keep the design language honest: the substitute source represents the external system, while the shortcut still proves the no-copy access pattern described in the OneLake shortcut documentation \cite{microsoftOneLakeShortcuts}.
\end{notebox}

\subsection*{Reference Architecture}
\begin{figure}[H]
    \centering
    \resizebox{\textwidth}{!}{%
    \begin{tikzpicture}[
        capsource/.style={rectangle, rounded corners, draw=codegray, thick, fill=white, align=center, minimum width=2.9cm, minimum height=1.15cm, font=\small\bfseries},
        capvirtual/.style={rectangle, rounded corners, draw=primary, thick, fill=primary!8, align=center, minimum width=3.2cm, minimum height=1.15cm, font=\small\bfseries},
        caplake/.style={cylinder, shape border rotate=90, aspect=0.25, draw=primary, thick, fill=backcolour, align=center, minimum width=2.8cm, minimum height=1.45cm, font=\small\bfseries},
        capconsume/.style={rectangle, rounded corners, draw=codegreen!60!black, thick, fill=green!7, align=center, minimum width=2.7cm, minimum height=1.1cm, font=\small\bfseries},
        caparrow/.style={draw, thick, -Latex, draw=primary}
    ]
        \node[capsource] (s3cap) at (0, 1.45) {AWS S3\\Historical Sales};
        \node[capsource] (dvcap) at (0, -1.45) {Dataverse\\Current CRM};
        \node[capvirtual] (shortcutcap) at (4.1, 0) {OneLake Shortcuts\\Read in Place};
        \node[caplake, fill=brown!12] (bronzecap) at (8.0, 0) {Lakehouse Bronze\\Virtual Raw};
        \node[capvirtual, fill=blue!6] (sparkcap) at (11.9, 0) {PySpark Notebook\\Standardize + Join};
        \node[caplake, fill=gray!15] (silvercap) at (15.8, 0) {Silver Delta\\Curated Table};
        \node[capconsume] (pbicap) at (19.5, 0) {Power BI\\Direct Lake Ready};

        \draw[caparrow] (s3cap) -- (shortcutcap);
        \draw[caparrow] (dvcap) -- (shortcutcap);
        \draw[caparrow] (shortcutcap) -- (bronzecap);
        \draw[caparrow] (bronzecap) -- (sparkcap);
        \draw[caparrow] (sparkcap) -- (silvercap);
        \draw[caparrow] (silvercap) -- (pbicap);

        \node[note, text width=13.5cm, align=center] at (9.7, -2.35) {Raw source data remains in AWS or Dataverse. Fabric materializes only the governed Silver Delta-Parquet table that downstream consumers can trust.};
    \end{tikzpicture}%
    }
    \caption{Capstone architecture for a multi-cloud retail landing zone using OneLake shortcuts.}
    \label{fig:chapter1-capstone-retail-landing-zone}
\end{figure}

\subsection*{Implementation Steps}
\begin{enumerate}
    \item \textbf{Create the workspace.} In Fabric, create \texttt{fabric-de-week01-retail}. Assign trial capacity or the instructor-approved capacity. Record the capacity name and whether the workspace is intended for development only.
    \item \textbf{Create the Lakehouse.} Add \texttt{lh\_retail\_landing}. Confirm that the Lakehouse exposes \texttt{Files} and \texttt{Tables}; the Silver output will become a managed Delta table, not a loose folder of files \cite{microsoftFabricLakehouse}.
    \item \textbf{Create the Bronze folders.} Under \texttt{Files}, use \texttt{bronze/s3\_sales} for historical sales and \texttt{bronze/dataverse\_crm} for current CRM. These names describe source origin and medallion intent.
    \item \textbf{Create the sales shortcut.} Add a OneLake shortcut at \texttt{Files/bronze/s3\_sales}. Use Amazon S3 when credentials are available. Otherwise, point the shortcut to a public Blob folder that simulates historical sales. Do not upload the sales files manually.
    \item \textbf{Create the CRM shortcut.} Add a shortcut at \texttt{Files/bronze/dataverse\_crm}. Use Dataverse when permitted. Otherwise, point to simulated CRM profile and loyalty files.
    \item \textbf{Validate virtualization.} Open each shortcut folder and confirm that files are visible through OneLake. Document the remote source, authentication method, owner, and expected refresh pattern.
    \item \textbf{Create the notebook.} Add \texttt{nb\_capstone\_retail\_silver}, attach it to \texttt{lh\_retail\_landing}, and run a small row-count profile before transformation.
    \item \textbf{Write the Silver table.} Use the notebook pattern in \cref{lst:chapter1-capstone-pyspark} and adjust column names to match your source files.
\end{enumerate}

\begin{lstlisting}[language=Python,caption={PySpark notebook pattern for reading shortcut data in place and writing a Silver Delta table.},label={lst:chapter1-capstone-pyspark}]
from pyspark.sql import functions as F

sales_path = "Files/bronze/s3_sales/**/*.parquet"
crm_path = "Files/bronze/dataverse_crm/**/*.parquet"

sales_raw_df = spark.read.parquet(sales_path)
crm_raw_df = spark.read.parquet(crm_path)

sales_df = (
    sales_raw_df
    .select(
        F.col("OrderId").cast("string").alias("order_id"),
        F.col("CustomerId").cast("string").alias("customer_id"),
        F.to_date("OrderDate").alias("order_date"),
        F.col("Region").cast("string").alias("region"),
        F.col("SalesAmount").cast("decimal(18,2)").alias("sales_amount"),
        F.col("CurrencyCode").cast("string").alias("currency_code"),
    )
    .where(F.col("order_id").isNotNull())
    .dropDuplicates(["order_id"])
)

crm_df = (
    crm_raw_df
    .select(
        F.col("CustomerId").cast("string").alias("customer_id"),
        F.col("LoyaltyTier").cast("string").alias("loyalty_tier"),
        F.col("EmailOptIn").cast("boolean").alias("email_opt_in"),
        F.to_timestamp("ModifiedOn").alias("crm_modified_on"),
    )
    .where(F.col("customer_id").isNotNull())
    .dropDuplicates(["customer_id"])
)

silver_df = (
    sales_df
    .join(crm_df, on="customer_id", how="left")
    .withColumn("sales_year", F.year("order_date"))
    .withColumn("silver_loaded_at", F.current_timestamp())
)

silver_df.write.format("delta").mode("overwrite").saveAsTable(
    "silver_retail_customer_sales"
)
\end{lstlisting}

The second listing is a documentation artifact you can store with your runbook or deployment notes. It records the intended shortcut layout without placing secrets in source control. Real deployments should use secure connection objects, tenant-approved credentials, and environment-specific automation.

\begin{lstlisting}[language=json,caption={Shortcut intent record for the capstone landing zone.},label={lst:chapter1-capstone-shortcut-json}]
{
  "workspace": "fabric-de-week01-retail",
  "lakehouse": "lh_retail_landing",
  "bronzeShortcuts": [
    {
      "name": "s3_sales",
      "oneLakePath": "Files/bronze/s3_sales",
      "sourceType": "AmazonS3",
      "sourceUri": "s3://retail-history-example/sales/",
      "copyData": false,
      "owner": "retail-platform-history"
    },
    {
      "name": "dataverse_crm",
      "oneLakePath": "Files/bronze/dataverse_crm",
      "sourceType": "Dataverse",
      "sourceUri": "environment://crm/customers",
      "copyData": false,
      "owner": "customer-experience-platform"
    }
  ],
  "silverOutput": {
    "table": "silver_retail_customer_sales",
    "format": "Delta",
    "materializedInOneLake": true
  }
}
\end{lstlisting}

If your class prefers SQL validation, run a SQL cell or SQL analytics endpoint query after the write:

\begin{lstlisting}[language=SQL,caption={SQL validation for the Silver customer sales table.}]
SELECT
    COUNT(*) AS row_count,
    COUNT(DISTINCT customer_id) AS customer_count,
    MIN(order_date) AS first_order_date,
    MAX(order_date) AS last_order_date,
    SUM(sales_amount) AS total_sales_amount
FROM silver_retail_customer_sales;
\end{lstlisting}

\subsection*{Deliverables}
Submit a concise package that proves the architecture and implementation:
\begin{itemize}
    \item Workspace name, capacity assignment, Lakehouse name, and notebook name.
    \item Screenshot or exported evidence showing both shortcut folders under \texttt{Files/bronze}.
    \item A short runbook entry for each shortcut: source owner, source URI pattern, authentication method, update frequency, and escalation contact.
    \item Notebook output showing row counts for historical sales, CRM profiles, and the final Silver table.
    \item The schema of \texttt{silver\_retail\_customer\_sales}, including keys, dates, measures, and CRM enrichment fields.
    \item A paragraph explaining which data was virtualized, which data was materialized, and why that split is appropriate for Week~1.
\end{itemize}

\subsection*{Acceptance Criteria}
Use this rubric before declaring the capstone complete:
\begin{itemize}
    \item[\(\square\)] The workspace exists on trial or approved Fabric capacity and contains the required Lakehouse.
    \item[\(\square\)] Both external sources are represented through OneLake shortcuts under the Bronze zone.
    \item[\(\square\)] Historical sales data is read in place and not bulk copied into OneLake.
    \item[\(\square\)] CRM or simulated CRM data is read in place and not bulk copied into OneLake.
    \item[\(\square\)] The PySpark notebook reads both shortcut paths successfully and includes basic row-count or null-count validation.
    \item[\(\square\)] The notebook writes \texttt{silver\_retail\_customer\_sales} as a managed Delta table.
    \item[\(\square\)] The Silver table has standardized customer keys, dates, currency or amount fields, and CRM enrichment fields.
    \item[\(\square\)] The design explicitly distinguishes virtual Bronze access from materialized Silver curation.
    \item[\(\square\)] The final explanation states that raw source data remains in the external systems and that only curated output is stored as Delta-Parquet in the Lakehouse.
    \item[\(\square\)] No secrets, connection strings, access keys, or personal data are pasted into notebook markdown, listings, screenshots, or source-controlled files.
\end{itemize}

\subsection*{Stretch Goals}
If you finish early, extend the capstone without changing the core no-copy promise:
\begin{itemize}
    \item Add a Gold aggregate table such as \texttt{gold\_regional\_loyalty\_sales} grouped by region, sales year, and loyalty tier.
    \item Build a small Power BI semantic model over the Silver or Gold table and test whether Direct Lake is appropriate for the table shape.
    \item Add an additional shortcut to ADLS Gen2 to represent a supplier, inventory, or product catalog feed.
    \item Convert the notebook into a repeatable medallion job with parameters for source paths, target table name, and write mode.
    \item Add data quality assertions for duplicate order identifiers, invalid dates, negative sales amounts, missing customer keys, and unsupported currency codes.
    \item Document a promotion plan from development to test to production workspaces, including who owns shortcut credentials and who certifies the Gold semantic model.
\end{itemize}
